% ---------------------------------------------------------------------------
% NasAge: main figures 5 and 6, with legends.
%
% This file is self-contained and makes NO reference to supplementary figures.
% The supplementary set is built separately from figures_supplementary.tex.
%
% Layout: each figure occupies a full page, scaled to the text block, with its
% legend on the page that follows. The figures are 180 x 202-252 mm, so a figure
% and its legend cannot share an A4 page without shrinking the panels below
% legibility.
%
% Compile with:  tectonic -X compile figures_main.tex
%            or: pdflatex figures_main.tex
% The figure PDFs must sit beside this file.
% ---------------------------------------------------------------------------
\documentclass[11pt,a4paper]{article}

\usepackage[margin=20mm]{geometry}
\usepackage{graphicx}
\usepackage{amsmath}
\usepackage{microtype}
\usepackage[hidelinks]{hyperref}

\setlength{\parindent}{0pt}
\setlength{\parskip}{6pt}

\newcommand{\figpage}[1]{%
  \clearpage
  \thispagestyle{empty}%
  \vspace*{\fill}
  \begin{center}
    \includegraphics[width=\textwidth,height=\textheight,keepaspectratio]{#1}
  \end{center}
  \vspace*{\fill}
}

\newcommand{\legend}[3]{%
  \clearpage
  {\large\textbf{Figure #1}\par}
  \vspace{2pt}
  {\textbf{#2}\par}
  \vspace{6pt}
  \small #3\par
}

\title{NasAge: Figures 5 and 6}
\author{}
\date{Generated 25 November 2026}

\begin{document}
\maketitle

\section*{Why these two figures exist}

Figures 1 to 4 establish an age-associated remodelling of the nasal epithelium, derived
within our own cohort.

\textbf{Figure 1} shows that basal cells decrease in abundance with age, and that a
pre-trained transcriptomic clock (Pasta; Salignon \emph{et al.}, \emph{Advanced Science},
\texttt{doi:10.1002/advs.76740}) tracks donor age in the epithelium, establishing that the
tissue carries a readable age signal.
\textbf{Figure 2} resolves that loss: the basal compartment is not uniform. Basal 1 does not
change with age, whereas Basal 2 advances along pseudotime with age, so the change is a
shift in differentiation state rather than a simple depletion.
\textbf{Figure 3} shows that senescence is cell-type-restricted. CDKN2A$^{+}$ ciliated cells
accumulate with age and show reduced cilium movement together with increased senescence and
SASP programmes; both the decrease in basal cells and the progression of basal pseudotime
correlate with the rise in p16$^{+}$ ciliated cells, linking the two compartments.
\textbf{Figure 4} turns to immune tone. Interferon programmes peak in childhood, in line with
Yoshida \emph{et al.}, \emph{Nature} 2022, and the response is bimodal: classical interferon-
stimulated genes peak only in childhood, while a distinct interferon-signalling gene cluster
is high again in the elderly.

Together these four figures describe a coherent trajectory in a single cohort. That is also
their limitation, and it defines what the next two figures must do.

\textbf{Figure 5 asks whether the signal is real outside the cohort that produced it.} Every
claim above rests on one dataset, analysed by us, with analytic choices made while looking at
it. Figure 5 therefore re-tests the mucociliary and interferon components in independent
nasal cohorts under plans fixed in advance, and adds the control that a correlation with age
most often fails: whether randomly chosen gene sets of matched expression and size behave the
same way.

\textbf{Figure 6 asks whether these programmes are a usable measurement, and whether they
matter.} Figures 1 to 4 treat gene programmes as readouts of epithelial state, which is only
legitimate if each programme reports on the cell type it is meant to report on. Figure 6
tests that directly across three cohorts, re-examines the age trajectory of the one immune
axis that survives, and then asks the question the preceding figures cannot: whether
epithelial state measured before exposure relates to what happens on challenge with a
pathogen.

\section*{How these figures were produced}

Every panel is drawn from a results file generated by an analysis plan committed to version
control \emph{before} the corresponding data were accessed. Each plan fixed the cohorts, the
unit of analysis, the covariates, the direction of effect, the multiple-testing family and
the verdict table; deviations are dated in the results file rather than edited into the plan.
No panel recomputes a model statistic, and the plotting scripts only read results and arrange
them.

\textbf{Data and independence.} Validation cohorts are public nasal transcriptomic datasets
held outside the analysis repository. No dataset that contributed to the primary findings in
Figures 1 to 4 is used here, and the Yoshida 2022 nasal atlas, together with any integrated
collection containing its donors, is excluded from every panel.

\textbf{Scoring and units.} Gene programmes are scored per cell and aggregated to donor
$\times$ cell-type pseudobulk, requiring at least 20 cells per unit, or scored directly on
bulk samples. A programme is tested only if it passes a measurability gate of at least 60\%
of its genes detected in at least 5\% of cells. Because the choice of scoring statistic is
itself an analytic degree of freedom, four are carried through in parallel: mean $z$,
median $z$, a rank statistic computed against the whole transcriptome, and the first
principal component of the set.

\textbf{Models.} Age effects are linear models on donor $\times$ cell-type units with
sequencing depth and sex as covariates. Infection outcomes are logistic models on baseline
samples with a cell-composition proxy, combined across cohorts by fixed-effect
inverse-variance meta-analysis. Marker genes used to build composition proxies are excluded
from the programmes they adjust.

\textbf{Controls.} Two are applied throughout. Benjamini--Hochberg correction is taken over
the declared family, and the effective number of independent tests is reported beside it,
since correlated modules do not constitute as many tests as they are counted as. Separately,
each result is compared against 500 random gene sets matched to the tested set on expression
level and size, scored by the identical pipeline; the reported empirical $p$ is the
\emph{worst} across the four scoring statistics, a rule fixed before any null was drawn.
Where an association is central, a specification curve reports every combination of scoring
statistic and covariate set rather than a single preferred model.

\textbf{Rendering.} Figures are composed in Matplotlib at 180\,mm width with Type-42
embedded fonts and a colourblind-safe palette, and are written directly to vector PDF.

% ===========================================================================
\figpage{figure5.pdf}

\legend{5}{The mucociliary ageing signal replicates in an independent nasal cohort.}{%
\textbf{(a)} M02 mucociliary axoneme score against donor age in the replication cohort
($n=120$ donors, bulk nasal brushing). Line and band are the ordinary least-squares fit and
its 95\% confidence interval. The Spearman $\rho$ is computed from the plotted points; the
adjusted $\beta$ is the age coefficient from the pre-registered model, which includes
sequencing depth as a covariate. Benjamini--Hochberg correction is over the declared family
of 13 modules, in which this test does \emph{not} survive ($\mathrm{BH}=0.22$).
\par
\textbf{(b)} The same analysis for CILIUM28, the gene set derived in the primary cohort.
Because CILIUM28 was defined on the primary data, this is a transfer test, not an independent
derivation, and it is reported alongside (a) rather than as confirmation of it.
\par
\textbf{(c)} Both age coefficients against the distribution of 500 expression- and
size-matched random gene sets scored identically in the same cohort. Neither observed
coefficient is approached by any random set, and \textbf{0 of 500} random sets reach $p<0.05$
against a nominal 5\%. The association is therefore specific to these gene sets rather than a
property of any set of comparably expressed genes.
\par
\textbf{(d)} Interferon-stimulated gene score (M05) in children versus adults in an
independent paediatric and adult cohort, by epithelial cell type, as donor means; horizontal
bars are group medians. Benjamini--Hochberg values are within the declared family, and
\emph{emp.\ p} is the expression-matched empirical $p$, with an asterisk marking the
combinations that clear it. The ciliated contrast reaches $\mathrm{BH}=1\times10^{-3}$ but
fails its matched null ($\mathrm{emp.}\ p=0.076$) and is not claimed.
\par
\textbf{(e)} The same contrast for M06 MHC-I processing, which clears both the family-wise
correction and the matched null in all three cell types.
\par
\textbf{Interpretation.} Two different questions are being asked here and they give different
answers. The matched null in (c) asks whether \emph{these} gene sets behave unlike comparably
expressed random sets, and the answer is unambiguous: 0 of 500. The family-wise correction
asks whether this test stands out among the 13 modules examined, and there the M02 result
does not survive ($\mathrm{BH}=0.22$). Both are reported because neither alone is the honest
summary: the signal is specific to the gene set but modest relative to the family it was
tested in. Panels (d) and (e) then locate where in life the signal sits. It is a
childhood-to-adult step rather than a gradient, since within adulthood the slope is flat in
two further independent cohorts. The one combination that reaches family-wise significance
but fails its matched null, ciliated M05, is the useful negative here, because it shows the
matched null is doing real work rather than rubber-stamping whatever the $p$-value says.
\par
\textbf{Conclusion.} The mucociliary and interferon ageing signal of Figures 1 and 4
reproduces in independent nasal cohorts and is specific to the curated gene sets rather than
to expression level. Its age structure is a childhood-to-adult difference, so it should be
described as a maturational contrast and not as a continuous ageing clock. The M02 axoneme
replication is supported as a specific effect but not as a family-wise significant one, and
should be reported with both figures attached.}

% ===========================================================================
\figpage{figure6.pdf}

\legend{6}{A cell-type-specific mucosal programme panel, and its relation to infection.}{%
\textbf{(a)} Mean score of each curated programme in each epithelial family, across three
single-cell cohorts and 224 donor $\times$ cell-type units. The solid box marks the cell type
each programme was curated for; the dashed orange box marks the one programme (P6 viral entry
factors) that scores highest somewhere else.
\par
\textbf{(b)} Standardised mean difference between a programme's own cell type and the other
two, with 95\% confidence intervals. Five of six programmes are specific and the sign is
unanimous in every cohort. P6 is significantly \emph{lower} in its assigned cell type
($\mathrm{SMD}=-1.66$) and is reported as misassigned rather than dropped. Curated cell-type
genes are expected to be enriched in that cell type, so the magnitude of these differences
overstates the achievement; the informative result is the one programme that fails.
\par
\textbf{(c)} M06 antigen presentation against age, 4--77 years, as donor $\times$ cell-type
means with per-family least-squares fits. The decline is present in all three cell types
($\mathrm{BH}\le1\times10^{-4}$), is robust across 12 analytic specifications, and clears its
expression-matched null under all four scoring methods in basal and secretory cells (three of
four in ciliated). As in Figure 5, the slope is carried by the childhood-to-adult contrast;
within adulthood it is flat.
\par
\textbf{(d)} Odds ratio per standard deviation of baseline P5 mucin barrier score for
infection outcome, in three challenge cohorts, with a fixed-effect inverse-variance pooled
estimate. The direction is protective in all three. Only pneumococcal carriage clears its own
expression-matched null (f); RSV and influenza do not, so the pooled estimate is shown and
not claimed.
\par
\textbf{(e)} Specification curve for P5 in the RSV and pneumococcal cohorts: 16
specifications each, being four scoring methods $\times$ four covariate sets, ordered by
estimate; filled points are $p<0.05$. All 32 specifications fall below $\mathrm{OR}=1$, that
is 100\% sign stability in both cohorts. Stability of this kind is evidence that a
multiplicity correction can neither supply nor remove.
\par
\textbf{(f)} The observed P5 effect in pneumococcal carriage against 500 expression- and
size-matched random gene sets scored identically. The empirical $p$ is 0.000 under every one
of the four scoring methods. By a rule fixed before any null was drawn, the figure reported
for each cohort is the \emph{worst} scorer, never the best.
\par
\textbf{Interpretation.} The figure separates into three claims of decreasing strength.
Panels (a) and (b) establish the panel as a \emph{measurement instrument}: the programmes
read out the cell types they were built for, with one exception that is reported rather than
removed. This is the precondition for reading Figures 1 to 4 as statements about epithelial
state at all. Panel (c) is a reproducible biological association and the strongest trajectory
result we have. Panels (d) to (f) are the weakest and most interesting: the direction of the
P5--infection association is completely stable, surviving all 32 analytic specifications in
two cohorts, yet stability and specificity are different properties and only pneumococcal
carriage has both. The natural reading of (e) and (f) together is that P5 captures something
real about susceptibility in at least one setting, and that the two other cohorts are
consistent with it but cannot independently support it. Panel (f) is also where the reporting
rule matters most: empirical $p$ for P5 ranges from 0.032 to 0.694 across scorers in the
influenza cohort, so a figure chosen after the fact could have been made to say almost
anything. Fixing the worst-scorer rule in advance is what makes the pneumococcal result,
0.000 under all four, credible.
\par
\textbf{Limits.} We could obtain no vaccine outcome in any reachable cohort. The one vaccine
exposure that could be tested, a live attenuated intranasal influenza vaccine, did not engage
any of these programmes despite a positive control that rose as expected in the same samples,
confirming the assay works. P5's outcome association does not survive if all 140 primary
tests across this programme of work are treated as a single family. Challenge cohorts are
small, at 39, 39 and 74 participants. The cell-type specificity in (a) and (b) is partly
definitional, since curated marker genes should be enriched where they were curated.
\par
\textbf{Conclusion.} The programme panel is a valid cell-type-resolved measurement of
epithelial state, which supports the use of programme scores throughout this work, and M06
antigen presentation carries a reproducible, specificity-tested decline with age. The P5
mucin barrier axis is a stable correlate of reduced infection susceptibility, demonstrated as
specific in pneumococcal carriage and consistent but unproven in RSV and influenza challenge.
\textbf{This figure must not be presented as a mucosal vaccine readiness index}: no vaccine
outcome was available in any reachable cohort, and the one vaccine exposure that could be
tested produced a null against a working control.}

\end{document}
